\section{Análisis de riesgos}

El análisis de riesgos se realiza tomando en cuenta el tiempo estimado de desarrollo de aproximadamente 6 meses, los recursos disponibles para el desarrollo tecnológico y las habilidades técnicas para el diseño de simuladores de autómatas celulares, cómputo paralelo, redes neuronales y programación orientada a eventos. 

\subsection{Identificación de riesgos}

Los riesgos se agrupan en cinco categorías según su origen: técnicos, operativos, de requerimientos, de estimación y de alcance. A cada riesgo se le asigna un código formado por las iniciales de su categoría y un número consecutivo, el cual se utiliza en el resto del análisis.

\subsubsection{Riesgos técnicos}

\begin{enumerate}[label*=RT\arabic*, noitemsep]
	\item Incompatibilidad entre diferentes lenguajes de programación: Los módulos desarrollados en Python, C++ y CUDA generan errores de integración al intercambiar información y sincronizar procesos.
	
	\item Indisponibilidad de \textit{hardware}: La tarjeta gráfica, el procesador u otros componentes no tienen el rendimiento suficiente.
	
	\item Incompatibilidad entre tecnologías: El sistema operativo, los lenguajes de programación y los elementos de \textit{hardware} no son compatibles para el desarrollo y ejecución.
	
	\item Dependencias de bibliotecas: La instalación y configuración de diferentes bibliotecas varía según el sistema operativo.
	
	\item Mala visualización: El renderizado de los componentes gráficos es incorrecto o está mal optimizado.
	
	\item Indisponibilidad de \textit{software} con licencia: Las herramientas que requieren una licencia, como Wolfram Mathematica, dejan de estar disponibles durante el desarrollo por la cancelación o el vencimiento de la licencia.
\end{enumerate}

\subsubsection{Riesgos operativos}
\begin{enumerate}[label*=RO\arabic*, noitemsep]
	\item Documentación deficiente: La documentación no permite mantener ni ampliar el proyecto a futuro.
	
	\item Problemas personales: Falta de motivación o problemas de salud física o mental se presentan durante el desarrollo.
	
	\item Curva de aprendizaje elevada: Se encuentra una curva de aprendizaje elevada y se carece de experiencia con la tecnología a utilizar.
\end{enumerate}
\subsubsection{Riesgos de requerimientos}
\begin{enumerate}[label*=RR\arabic*, noitemsep]
	\item Requerimientos ambiguos: La definición incompleta o ambigua de los requerimientos puede provocar cambios en etapas avanzadas del desarrollo.
\end{enumerate}

\subsubsection{Riesgos de estimación}
\begin{enumerate}[label*=RE\arabic*, noitemsep]
	\item Subestimación del tiempo de desarrollo: La complejidad de enlazar diferentes tecnologías puede extender los plazos planeados.
\end{enumerate}

\subsubsection{Riesgos de alcance}
\begin{enumerate}[label*=RA\arabic*, noitemsep]
	\item Limitación del alcance funcional: Las restricciones de tiempo pueden dar pie a limitar el número de funcionalidades requeridas.

	\item Carga de trabajo imprevista: El poder de cómputo resulta insuficiente, por lo que se reduce la complejidad de las simulaciones.
\end{enumerate}

\subsection{Probabilidad e impacto}

El nivel de cada riesgo se obtiene a partir de dos factores. La probabilidad es la posibilidad de que el riesgo ocurra durante el desarrollo y se clasifica en cinco niveles: muy poco probable (10\%), poco probable (30\%), moderadamente probable (50\%), probable (70\%) y muy probable (90\%). El impacto, o severidad, es la magnitud del daño que causaría al proyecto en caso de ocurrir y se clasifica en nulo, menor, moderado, crítico y fatal. La matriz de riesgos de la tabla \ref{table:matriz_riesgos} combina ambos factores y asigna a cada riesgo un nivel bajo, medio, alto o muy alto \cite{matriz_riesgos}. Siguiendo dicha tabla, en las tablas \ref{table:riesgos_tecnicos} a \ref{table:riesgos_alcance} se enlistan los niveles de riesgo para cada riesgo identificado.

\begin{table}[H]
	\centering
	\renewcommand{\arraystretch}{1.2}
	\setlength{\tabcolsep}{3pt}
	\begin{tabular}{|>{\centering\arraybackslash}m{0.04\textwidth}|p{0.13\textwidth}|*{5}{>{\centering\arraybackslash}p{0.14\textwidth}|}}
		\hhline{~~|-----|}
		\multicolumn{2}{c|}{} & \multicolumn{5}{c|}{\textbf{Probabilidad}} \\ \hhline{~~|-----|}
		\multicolumn{2}{c|}{} & Muy poco probable $10\%$ & Poco probable $30\%$ &  Moderada-mente $50\%$ & Probable $70\%$ & Muy probable $90\%$ \\ \hhline{|-------|}
		\multirow{5}{*}{\rotatebox[origin=c]{90}{\textbf{Severidad}}}
		& Nula & \cellcolor{riskLow}Bajo & \cellcolor{riskLow}Bajo & \cellcolor{riskLow}Bajo & \cellcolor{riskMedium}Medio & \cellcolor{riskMedium}Medio \\ \hhline{|~|------|}
		& Menor & \cellcolor{riskLow}Bajo & \cellcolor{riskLow}Bajo & \cellcolor{riskMedium}Medio & \cellcolor{riskMedium}Medio & \cellcolor{riskMedium}Medio \\ \hhline{|~|------|}
		& Moderado & \cellcolor{riskMedium}Medio & \cellcolor{riskMedium}Medio & \cellcolor{riskMedium}Medio & \cellcolor{riskHigh}Alto & \cellcolor{riskHigh}Alto \\ \hhline{|~|------|}
		& Crítico & \cellcolor{riskMedium}Medio & \cellcolor{riskMedium}Medio & \cellcolor{riskHigh}Alto & \cellcolor{riskHigh}Alto & \cellcolor{riskVeryHigh}Muy alto \\ \hhline{|~|------|}
		& Fatal & \cellcolor{riskMedium}Medio & \cellcolor{riskHigh}Alto & \cellcolor{riskHigh}Alto & \cellcolor{riskVeryHigh}Muy alto & \cellcolor{riskVeryHigh}Muy alto \\ \hhline{|-------|}
	\end{tabular}
	\caption{Matriz de riesgos \cite{matriz_riesgos}}
	\label{table:matriz_riesgos}
\end{table}

\subsection{Plan de mitigación}

El plan de mitigación define, para cada riesgo identificado, las acciones preventivas que se realizan durante el desarrollo para reducir la probabilidad de que ocurra o su impacto. Las tablas \ref{table:riesgos_tecnicos} a \ref{table:riesgos_alcance} presentan, por categoría, el código y el nombre de cada riesgo junto con su plan de mitigación.

\begingroup
\renewcommand{\arraystretch}{1.2}
\setlength{\tabcolsep}{2.5pt}
\begin{longtable}{|>{\raggedright\arraybackslash}p{0.13\textwidth}|p{0.125\textwidth}|p{0.19\textwidth}|>{\raggedright\arraybackslash}p{0.26\textwidth}|p{0.105\textwidth}|p{0.11\textwidth}|}
	\endfirsthead
	\hline
	\endhead
		\hline
		\textbf{RT1} & \riskSpan{Incompatibilidad entre diferentes lenguajes de programación} \\* \hhline{|------|}\noalign{\nobreak}
		\textbf{Impacto:} & Crítico & \textbf{Probabilidad:} & Moderadamente & \textbf{Riesgo:} & \cellcolor{riskHigh}Alto \\* \hhline{|------|}\noalign{\nobreak}
		\textbf{Plan de mitigación:} & \riskSpan{Verificar la documentación de compatibilidad entre los lenguajes de programación y realizar pruebas de transferencia de datos y sincronización previas al desarrollo.} \\* \noalign{\hrule height 1pt}\noalign{\penalty0}
		\textbf{RT2} & \riskSpan{Indisponibilidad de \textit{hardware}} \\* \hhline{|------|}\noalign{\nobreak}
		\textbf{Impacto:} & Moderado & \textbf{Probabilidad:} & Poco probable & \textbf{Riesgo:} & \cellcolor{riskMedium}Medio \\* \hhline{|------|}\noalign{\nobreak}
		\textbf{Plan de mitigación:} & \riskSpan{Asegurar la disponibilidad de equipos de respaldo o entornos virtuales alternativos para continuar el desarrollo en caso de fallos de \textit{hardware}.} \\* \noalign{\hrule height 1pt}\noalign{\penalty0}
		\textbf{RT3} & \riskSpan{Incompatibilidad entre tecnologías} \\* \hhline{|------|}\noalign{\nobreak}
		\textbf{Impacto:} & Fatal & \textbf{Probabilidad:} & Probable & \textbf{Riesgo:} & \cellcolor{riskVeryHigh}Muy alto \\* \hhline{|------|}\noalign{\nobreak}
		\textbf{Plan de mitigación:} & \riskSpan{Realizar pruebas de compatibilidad del \textit{hardware} y \textit{software} en diferentes sistemas operativos y versiones antes de la integración.} \\* \noalign{\hrule height 1pt}\noalign{\penalty0}
		\textbf{RT4} & \riskSpan{Dependencias de bibliotecas} \\* \hhline{|------|}\noalign{\nobreak}
		\textbf{Impacto:} & Menor & \textbf{Probabilidad:} & Poco probable & \textbf{Riesgo:} & \cellcolor{riskLow}Bajo \\* \hhline{|------|}\noalign{\nobreak}
		\textbf{Plan de mitigación:} & \riskSpan{Documentar y verificar las versiones de bibliotecas utilizadas, además de considerar alternativas equivalentes que garanticen continuidad en caso de incompatibilidad.} \\* \noalign{\hrule height 1pt}\noalign{\penalty0}
		\textbf{RT5} & \riskSpan{Mala visualización} \\* \hhline{|------|}\noalign{\nobreak}
		\textbf{Impacto:} & Menor & \textbf{Probabilidad:} & Moderadamente & \textbf{Riesgo:} & \cellcolor{riskMedium}Medio \\* \hhline{|------|}\noalign{\nobreak}
		\textbf{Plan de mitigación:} & \riskSpan{Realizar pruebas de usabilidad y ajustar la interfaz de acuerdo con las resoluciones de pantalla y dispositivos objetivo.} \\* \noalign{\hrule height 1pt}\noalign{\penalty0}
		\textbf{RT6} & \riskSpan{Indisponibilidad de \textit{software} con licencia} \\* \hhline{|------|}\noalign{\nobreak}
		\textbf{Impacto:} & Crítico & \textbf{Probabilidad:} & Muy probable & \textbf{Riesgo:} & \cellcolor{riskVeryHigh}Muy alto \\* \hhline{|------|}\noalign{\nobreak}
		\textbf{Plan de mitigación:} & \riskSpan{Verificar la vigencia de las licencias al inicio del proyecto, limitar el uso de las herramientas con licencia a tareas que no formen parte del simulador y mantener alternativas de \textit{software} libre, como Python, para las tareas de análisis.} \\ \hline
	\caption{Nivel y plan de mitigación de los riesgos técnicos}
	\label{table:riesgos_tecnicos}
\end{longtable}
\endgroup

\begingroup
\renewcommand{\arraystretch}{1.2}
\setlength{\tabcolsep}{2.5pt}
\begin{longtable}{|>{\raggedright\arraybackslash}p{0.13\textwidth}|p{0.125\textwidth}|p{0.19\textwidth}|>{\raggedright\arraybackslash}p{0.26\textwidth}|p{0.105\textwidth}|p{0.11\textwidth}|}
	\endfirsthead
	\hline
	\endhead
		\hline
		\textbf{RO1} & \riskSpan{Documentación deficiente} \\* \hhline{|------|}\noalign{\nobreak}
		\textbf{Impacto:} & Moderado & \textbf{Probabilidad:} & Probable & \textbf{Riesgo:} & \cellcolor{riskHigh}Alto \\* \hhline{|------|}\noalign{\nobreak}
		\textbf{Plan de mitigación:} & \riskSpan{Establecer un estándar de documentación y realizar revisiones periódicas del avance para garantizar la calidad y completitud de los registros técnicos.} \\* \noalign{\hrule height 1pt}\noalign{\penalty0}
		\textbf{RO2} & \riskSpan{Problemas personales} \\* \hhline{|------|}\noalign{\nobreak}
		\textbf{Impacto:} & Crítico & \textbf{Probabilidad:} & Moderadamente & \textbf{Riesgo:} & \cellcolor{riskHigh}Alto \\* \hhline{|------|}\noalign{\nobreak}
		\textbf{Plan de mitigación:} & \riskSpan{Definir roles de respaldo o distribuir responsabilidades para minimizar el impacto de ausencias o contratiempos personales.} \\* \noalign{\hrule height 1pt}\noalign{\penalty0}
		\textbf{RO3} & \riskSpan{Curva de aprendizaje elevada} \\* \hhline{|------|}\noalign{\nobreak}
		\textbf{Impacto:} & Moderado & \textbf{Probabilidad:} & Muy probable & \textbf{Riesgo:} & \cellcolor{riskHigh}Alto \\* \hhline{|------|}\noalign{\nobreak}
		\textbf{Plan de mitigación:} & \riskSpan{Programar capacitaciones iniciales y establecer un periodo de adaptación con apoyo técnico y materiales de referencia.} \\ \hline
	\caption{Nivel y plan de mitigación de los riesgos operativos}
	\label{table:riesgos_operativos}
\end{longtable}
\endgroup

\begingroup
\renewcommand{\arraystretch}{1.2}
\setlength{\tabcolsep}{2.5pt}
\begin{longtable}{|>{\raggedright\arraybackslash}p{0.13\textwidth}|p{0.125\textwidth}|p{0.19\textwidth}|>{\raggedright\arraybackslash}p{0.26\textwidth}|p{0.105\textwidth}|p{0.11\textwidth}|}
	\endfirsthead
	\hline
	\endhead
		\hline
		\textbf{RR1} & \riskSpan{Requerimientos ambiguos} \\* \hhline{|------|}\noalign{\nobreak}
		\textbf{Impacto:} & Menor & \textbf{Probabilidad:} & Poco probable & \textbf{Riesgo:} & \cellcolor{riskLow}Bajo \\* \hhline{|------|}\noalign{\nobreak}
		\textbf{Plan de mitigación:} & \riskSpan{Revisar periódicamente los requerimientos para aclararlos y validar el entendimiento del equipo de desarrollo.} \\ \hline
	\caption{Nivel y plan de mitigación de los riesgos de requerimientos}
	\label{table:riesgos_requerimientos}
\end{longtable}
\endgroup

\begingroup
\renewcommand{\arraystretch}{1.2}
\setlength{\tabcolsep}{2.5pt}
\begin{longtable}{|>{\raggedright\arraybackslash}p{0.13\textwidth}|p{0.125\textwidth}|p{0.19\textwidth}|>{\raggedright\arraybackslash}p{0.26\textwidth}|p{0.105\textwidth}|p{0.11\textwidth}|}
	\endfirsthead
	\hline
	\endhead
		\hline
		\textbf{RE1} & \riskSpan{Subestimación del tiempo de desarrollo} \\* \hhline{|------|}\noalign{\nobreak}
		\textbf{Impacto:} & Fatal & \textbf{Probabilidad:} & Probable & \textbf{Riesgo:} & \cellcolor{riskVeryHigh}Muy alto \\* \hhline{|------|}\noalign{\nobreak}
		\textbf{Plan de mitigación:} & \riskSpan{Dividir el proyecto en fases pequeñas y manejables, aplicando seguimiento semanal y reajuste de cronograma según los resultados obtenidos.} \\ \hline
	\caption{Nivel y plan de mitigación de los riesgos de estimación}
	\label{table:riesgos_estimacion}
\end{longtable}
\endgroup

\begingroup
\renewcommand{\arraystretch}{1.2}
\setlength{\tabcolsep}{2.5pt}
\begin{longtable}{|>{\raggedright\arraybackslash}p{0.13\textwidth}|p{0.125\textwidth}|p{0.19\textwidth}|>{\raggedright\arraybackslash}p{0.26\textwidth}|p{0.105\textwidth}|p{0.11\textwidth}|}
	\endfirsthead
	\hline
	\endhead
		\hline
		\textbf{RA1} & \riskSpan{Limitación del alcance funcional} \\* \hhline{|------|}\noalign{\nobreak}
		\textbf{Impacto:} & Moderado & \textbf{Probabilidad:} & Moderadamente & \textbf{Riesgo:} & \cellcolor{riskMedium}Medio \\* \hhline{|------|}\noalign{\nobreak}
		\textbf{Plan de mitigación:} & \riskSpan{Definir claramente los objetivos y alcances desde el inicio, con aprobación formal de los responsables del proyecto antes de iniciar nuevas funcionalidades.} \\* \noalign{\hrule height 1pt}\noalign{\penalty0}
		\textbf{RA2} & \riskSpan{Carga de trabajo imprevista} \\* \hhline{|------|}\noalign{\nobreak}
		\textbf{Impacto:} & Crítico & \textbf{Probabilidad:} & Probable & \textbf{Riesgo:} & \cellcolor{riskHigh}Alto \\* \hhline{|------|}\noalign{\nobreak}
		\textbf{Plan de mitigación:} & \riskSpan{Implementar revisiones periódicas del plan de trabajo y ajustar el alcance en función de los recursos disponibles y los resultados obtenidos.} \\ \hline
	\caption{Nivel y plan de mitigación de los riesgos de alcance}
	\label{table:riesgos_alcance}
\end{longtable}
\endgroup

\subsection{Plan de contingencia}

El plan de contingencia define las acciones que se ejecutan cuando un riesgo se materializa a pesar de las medidas de mitigación. Se elabora únicamente para los riesgos de nivel muy alto, que son los siguientes:
\begin{itemize}
	\item \textbf{RT3} Incompatibilidad entre tecnologías.
	\item \textbf{RT6} Indisponibilidad de \textit{software} con licencia.
	\item \textbf{RE1} Subestimación del tiempo de desarrollo. 
\end{itemize}

Para cada uno se indican su probabilidad, su impacto, el plan de evasión (acciones para evitar que el riesgo ocurra) y las acciones de contingencia (qué hacer cuando ocurre).

\subsubsection{RT3 Incompatibilidad entre tecnologías}

\begin{table}[H]
	\centering
	\renewcommand{\arraystretch}{1.2}
	\setlength{\tabcolsep}{2.5pt}
	\begin{tabular}{|>{\raggedright\arraybackslash}p{0.13\textwidth}|p{0.09\textwidth}|p{0.19\textwidth}|>{\raggedright\arraybackslash}p{0.295\textwidth}|p{0.105\textwidth}|p{0.11\textwidth}|}
		\hline
		\textbf{RT3} & \riskSpan{Incompatibilidad entre tecnologías} \\ \hhline{|------|}
		\textbf{Impacto:} & Fatal & \textbf{Probabilidad:} & Probable (70\%) & \textbf{Riesgo:} & \cellcolor{riskVeryHigh}Muy alto \\ \hhline{|------|}
		\textbf{Plan de evasión:} & \riskSpan{Realizar pruebas de integración tempranas entre las distintas tecnologías implicadas. Documentar las versiones, dependencias y configuraciones utilizadas para garantizar compatibilidad. Establecer un entorno de desarrollo unificado que minimice diferencias entre plataformas.} \\ \hhline{|------|}
		\textbf{Plan de contingencia:} & \riskSpan{Sustituir o actualizar los componentes que generen incompatibilidad. Implementar adaptadores o módulos intermedios que permitan la interoperabilidad entre sistemas. Reasignar recursos técnicos para resolver los conflictos de integración y evitar retrasos prolongados.} \\ \hline
	\end{tabular}
	\caption{Plan de contingencia para RT3 Incompatibilidad entre tecnologías}
	\label{table:contingencia_incompatibilidad}
\end{table}

\subsubsection{RT6 Indisponibilidad de software con licencia}

\begin{table}[H]
	\centering
	\renewcommand{\arraystretch}{1.2}
	\setlength{\tabcolsep}{2.5pt}
	\begin{tabular}{|>{\raggedright\arraybackslash}p{0.13\textwidth}|p{0.09\textwidth}|p{0.19\textwidth}|>{\raggedright\arraybackslash}p{0.295\textwidth}|p{0.105\textwidth}|p{0.11\textwidth}|}
		\hline
		\textbf{RT6} & \riskSpan{Indisponibilidad de \textit{software} con licencia} \\ \hhline{|------|}
		\textbf{Impacto:} & Crítico & \textbf{Probabilidad:} & Muy probable (90\%) & \textbf{Riesgo:} & \cellcolor{riskVeryHigh}Muy alto \\ \hhline{|------|}
		\textbf{Plan de evasión:} & \riskSpan{Verificar al inicio del proyecto la vigencia y las condiciones de renovación de cada licencia. Mantener el simulador independiente de las herramientas con licencia, de modo que estas solo se usen para el análisis de resultados. Preparar alternativas de \textit{software} libre (NumPy, Matplotlib y Scikit-learn) para las tareas de análisis.} \\ \hhline{|------|}
		\textbf{Plan de contingencia:} & \riskSpan{Adquirir una licencia individual por el tiempo restante del proyecto. Migrar los cuadernos de análisis a las alternativas de \textit{software} libre si no es posible adquirir la licencia. Reajustar el cronograma para absorber el tiempo de migración.} \\ \hhline{|------|}
	\end{tabular}
	\caption{Plan de contingencia para RT6 Indisponibilidad de \textit{software} con licencia}
	\label{table:contingencia_licencia}
\end{table}

\subsubsection{RE1 Subestimación del tiempo de desarrollo}

\begin{table}[H]
	\centering
	\renewcommand{\arraystretch}{1.2}
	\setlength{\tabcolsep}{2.5pt}
	\begin{tabular}{|>{\raggedright\arraybackslash}p{0.13\textwidth}|p{0.09\textwidth}|p{0.19\textwidth}|>{\raggedright\arraybackslash}p{0.295\textwidth}|p{0.105\textwidth}|p{0.11\textwidth}|}
		\hline
		\textbf{RE1} & \riskSpan{Subestimación del tiempo de desarrollo} \\ \hhline{|------|}
		\textbf{Impacto:} & Fatal & \textbf{Probabilidad:} & Probable (70\%) & \textbf{Riesgo:} & \cellcolor{riskVeryHigh}Muy alto \\ \hhline{|------|}
		\textbf{Plan de evasión:} & \riskSpan{Realizar estimaciones realistas basadas en experiencias previas y tareas medibles. Incluir márgenes de tiempo de reserva en el cronograma para absorber imprevistos. Efectuar reuniones de seguimiento semanales para detectar desviaciones tempranas.} \\ \hhline{|------|}
		\textbf{Plan de contingencia:} & \riskSpan{Reasignar recursos o personal adicional para recuperar el tiempo perdido. Reducir temporalmente el alcance funcional o posponer características no críticas. Reajustar el cronograma de entregas y comunicar los cambios al equipo y partes interesadas.} \\ \hline
	\end{tabular}
	\caption{Plan de contingencia para RE1 Subestimación del tiempo de desarrollo}
	\label{table:contingencia_subestimacion}
\end{table}

\section{Estudio de factibilidad}

El proyecto cuenta con los recursos de \textit{hardware} necesarios para llevar a cabo simulaciones de alto nivel computacional. El equipo de cómputo disponible permite implementar técnicas de paralelización, asegurando la capacidad de realizar múltiples inferencias en tiempos razonables. Adicionalmente, se dispone de herramientas tecnológicas de alto rendimiento que brindan la ventaja computacional requerida para ejecutar la simulación de manera eficiente.

En cuanto a la factibilidad económica, la mayor parte del \textit{software} utilizado es de libre acceso, incluyendo Python, C++, Visual Studio Code y TeXstudio. Se utiliza también Wolfram Mathematica, que requiere una licencia anual de 195 dólares, la cual se encontraba disponible de manera gratuita para estudiantes del IPN al inicio del proyecto. Al cancelarse dicha licencia se adquirió una licencia de estudiante por un año (ver el plan de contingencia del riesgo RT6), un costo que no compromete la viabilidad financiera del proyecto.

Desde el punto de vista operativo, el diseño de la interfaz facilita la definición de reglas y la ejecución de simulaciones, optimizando el flujo de trabajo del grupo de investigadores. Esto permite acelerar significativamente el proceso de experimentación y análisis de resultados, maximizando la eficiencia del equipo durante el desarrollo del proyecto.

Finalmente, se cumple la factibilidad temporal, ya que el cronograma de desarrollo estimado permite realizar todas las simulaciones necesarias para el análisis posterior, asegurando que el proyecto se pueda completar dentro del tiempo previsto.


